\begin{thebibliography}{99}
\bibitem{Haberman77} R.~Haberman. Nonlinear transition layers---the second
Painleve transcendent. \emph{Studies in Applied Mathematics} \textbf{57}
(1977), 247--270. \href{https://doi.org/10.1002/sapm1977573247}{doi:10.1002/sapm1977573247}.
\bibitem{Haberman79} R.~Haberman. Slowly varying jump and transition phenomena
associated with algebraic bifurcation problems. \emph{SIAM J. Appl. Math.}
\textbf{37} (1979), 69--106.
\bibitem{Suleimanov92} B.~I.~Suleimanov. A ``nonlinear'' generalization of
special functions of wave catastrophes described by double integrals.
\emph{Mathematical Notes} \textbf{52} (1992), 1146--1149.
\bibitem{Suleimanov94} B.~I.~Suleimanov. Influence of weak nonlinearity on the
high-frequency asymptotics in caustic rearrangements.
\emph{Theor. Math. Phys.} \textbf{98} (1994), 132--138.
\bibitem{KiselevSuleimanov99} O.~M.~Kiselev, B.~I.~Suleimanov.
The solution of the Painleve equations as special functions of catastrophes,
defined by a rejection in these equations of terms with derivative.
\href{https://arxiv.org/abs/solv-int/9902004}{arXiv:solv-int/9902004} (1999).
\bibitem{Kiselev01} O.~M.~Kiselev. Hard loss of stability in Painleve-2 equation.
\emph{J. Nonlinear Math. Phys.} \textbf{8} (2001), 65--95.
\href{https://doi.org/10.2991/jnmp.2001.8.1.8}{doi:10.2991/jnmp.2001.8.1.8}.
Preprint: \href{https://arxiv.org/abs/solv-int/9902007}{arXiv:solv-int/9902007} (1999).
\bibitem{KiselevGlebov03} O.~M.~Kiselev, S.~G.~Glebov.
An asymptotic solution slowly crossing the separatrix near a saddle--centre
bifurcation point. \emph{Nonlinearity} \textbf{16} (2003), 327--362.
\bibitem{KiselevGlebov07} O.~M.~Kiselev, S.~G.~Glebov.
The capture into parametric autoresonance. \emph{Nonlinear Dynamics}
\textbf{48} (2007), 217--230.
\bibitem{ItsNovokshenov86} A.~R.~Its, V.~Yu.~Novokshenov.
\emph{The Isomonodromic Deformation Method in the Theory of Painleve Equations}.
Lecture Notes in Mathematics 1191. Springer, 1986.
\bibitem{FIKN06} A.~S.~Fokas, A.~R.~Its, A.~A.~Kapaev, V.~Yu.~Novokshenov.
\emph{Painleve Transcendents: The Riemann--Hilbert Approach}.
Mathematical Surveys and Monographs 128. AMS, 2006.
\bibitem{OlverTrogdon12} S.~Olver, T.~Trogdon.
Nonlinear steepest descent and the numerical solution of Riemann--Hilbert problems.
\href{https://arxiv.org/abs/1205.5604}{arXiv:1205.5604} (2012).
\bibitem{Kiselev24} O.~M.~Kiselev. Integral Formulas for the Painleve-2 Transcendent.
\emph{Regular and Chaotic Dynamics} \textbf{29} (2024), 838--852.
\href{https://doi.org/10.1134/S1560354724560041}{doi:10.1134/S1560354724560041}.
\bibitem{Kiselev26} O.~M.~Kiselev.
Integral Representations for Solutions of the Linearized Fourth Painleve Equation.
\href{https://arxiv.org/abs/2609.12876}{arXiv:2609.12876} (2026).
\bibitem{JimboMiwa81} M.~Jimbo, T.~Miwa. Monodromy preserving deformation of
linear ordinary differential equations with rational coefficients. II.
\emph{Physica D} \textbf{2} (1981), 407--448.
\bibitem{DormandPrince80} J.~R.~Dormand, P.~J.~Prince.
A family of embedded Runge--Kutta formulae.
\emph{J. Comput. Appl. Math.} \textbf{6} (1980), 19--26.
\end{thebibliography}
